\newpage
\section{Analiza problemu}

\subsection{Zastosowanie listy dwukierunkowej}
Lista dwukierunkowa jest strukturą danych wykorzystywaną w programowaniu do przechowywania i organizowania elementów, które można przeglądać w obu kierunkach. Stosuje się ją między innymi w odtwarzaczach muzycznych, historii przeglądania stron internetowych, nawigacji między otwartymi dokumentami oraz w implementacji niektórych kolejek i list. Jej zaletą jest możliwość dynamicznego dodawania i usuwania elementów bez konieczności przesuwania pozostałych danych.

\subsection{Sposób działania algorytmu}
Lista dwukierunkowa składa się z węzłów, z których każdy przechowuje wartość oraz dwa wskaźniki: na poprzedni i następny element. Pierwszy węzeł nie posiada poprzednika, a ostatni nie posiada następnika. Dodawanie elementu polega na utworzeniu nowego węzła i odpowiednim ustawieniu wskaźników, natomiast usuwanie wymaga połączenia ze sobą sąsiednich węzłów. Przeglądanie listy odbywa się przez przechodzenie między węzłami, a czyszczenie polega na usunięciu wszystkich elementów i zwolnieniu zajmowanej pamięci.

\subsection{Przykład działania listy dwukierunkowej}
Przykładowo, lista może przechowować liczby 10, 20 i 30.

Początkowo lista jest pusta:
\begin{verbatim}
NULL
\end{verbatim}

Po dodaniu liczb 10, 20 i 30 na koniec listy otrzymujemy:
\begin{verbatim}
NULL <- 10 <-> 20 <-> 30 -> NULL
\end{verbatim}

Przechodząc od początku, program wyświetli liczby w kolejności: 10, 20, 30. Przechodząc od końca, wyświetli: 30, 20, 10.

Jeżeli usuniemy element o wartości 20, wskaźniki zostaną zaktualizowane, a lista będzie wyglądać następująco:
\begin{verbatim}
NULL <- 10 <-> 30 -> NULL
\end{verbatim}

W ten sposób można dodawać, usuwać i przeglądać liczby bez konieczności przesuwania pozostałych elementów.

\subsection{Opis wykorzystanych narzędzi}
\textbf{Git} to system kontroli wersji, który umożliwia zapisywanie historii zmian w projekcie, porównywanie wersji kodu oraz przywracanie wcześniejszych stanów programu.

\textbf{GitHub} to internetowa platforma służąca do przechowywania repozytoriów Git i współpracy nad projektami. Umożliwia przesyłanie zmian na serwer, pobieranie projektu do różnych lokalizacji oraz synchronizowanie pracy.

\textbf{Doxygen} to narzędzie służące do automatycznego generowania dokumentacji na podstawie komentarzy umieszczonych w kodzie źródłowym. Może tworzyć dokumentację w różnych formatach, w tym w formacie LaTeX, z którego można wygenerować dokument PDF.

\textbf{LaTeX} to system przygotowywania dokumentów, szczególnie przydatny przy tworzeniu dokumentacji technicznej i naukowej. W projekcie może zostać wykorzystany do opracowania dokumentacji projektowej oraz wygenerowania dokumentu PDF.